%% !TEX root = manuscript.tex

\chapter*{Conclusion}

Au terme de cette alternance, je peux revenir sur la question qui a guidé ce mémoire : comment
reconstituer et fiabiliser les données d'une entreprise au cours d'un changement d'ERP, afin
d'en tirer des outils de pilotage utilisables par les services métiers ?

Chacune de mes missions a apporté un élément de réponse. Le premier semestre a été celui de la
reconstruction du socle. La migration de Grimmo vers SPO ne consistait pas seulement à
transférer des données, mais à retrouver les règles d'un modèle qui n'étaient documentées nulle
part, à résoudre les problèmes d'intégrité référentielle que posaient l'ordre d'injection et les
identifiants dynamiques, et à adapter la technique aux impératifs du métier, avec la stratégie
des « coquilles vides » qui a permis aux équipes de facturer dès le basculement.

Le second semestre a prolongé ce travail, puis l'a valorisé. La reprise des tiers a achevé de
constituer le socle commercial du nouvel ERP, en créant 1~434 clients acquéreurs sans doublon
détecté par les trois clés d'unicité retenues. Le Copilote Financier a ensuite exploité ce socle
pour outiller la direction financière sur trois questions de gestion, de bout en bout, de la
compréhension des exports jusqu'à une application que les contrôleurs de gestion ouvrent
eux-mêmes. En parallèle, le maintien du système décisionnel et le déploiement du SSO ont assuré
la continuité de ce qui existait déjà.

C'est en comparant ces missions que se dégage ce que je retiens principalement de l'année. À
chaque fois, la question posée par un service supposait, pour y répondre, de construire d'abord
une catégorie qui n'existait pas dans les données : le rang d'une commande, l'identité d'un
client, la dérive d'une marge. Et à chaque fois, la qualité de cette construction a déterminé ce
que je pouvais produire ensuite. Le Copilote Financier en donne la démonstration la plus nette :
là où je disposais d'une profondeur temporelle, le modèle a produit un résultat exploitable ; là
où je n'avais qu'une photographie à date, il n'a pas pu en produire. Structurer une donnée et la
modéliser ne sont donc pas deux étapes indépendantes, et c'est la première qui décide de ce que
la seconde pourra en tirer.

Il reste beaucoup à faire. La migration n'est pas achevée, l'historique budgétaire reste à
reprendre, et le Copilote Financier ne deviendra un véritable outil de suivi que lorsque les
exports auront été historisés sur une durée suffisante. Ces chantiers ne sont pas des
imperfections du travail réalisé, ce sont les conditions de ce qui viendra après.

Au-delà des livrables, cette année a surtout fait évoluer ma manière de travailler. J'ai acquis
des réflexes que je considère comme durables : prouver plutôt que supposer, documenter ce que je
découvre, et adapter ma manière de vérifier à la nature de ce que je livre. J'ai également
mesuré à quel point, en entreprise, la communication compte autant que la technique : une
analyse rigoureuse n'a de valeur que si elle est comprise et acceptée par ceux à qui elle est
destinée.

Cette alternance restera une étape décisive de mon parcours. Elle a marqué mon passage du cadre
des études à celui de l'entreprise, avec ses exigences propres, et m'a permis de gagner en
autonomie comme en maturité professionnelle. J'en retiens autant des compétences techniques que
des enseignements humains, et la confirmation d'une voie qui me correspond. C'est donc avec
confiance, et l'envie de continuer à apprendre, que j'aborde la suite de mon chemin
professionnel.
